\begin{abwframe}
\node[anchor=north west,inner sep=0pt,outer sep=0pt] at (38.551,-150.803) {\begin{minipage}[t][363.546bp][t]{877.846bp}\setlength{\parindent}{0pt}\setlength{\parskip}{0pt}{\TimesNR\fontsize{24}{26.88}\selectfont\color{c000000}\itshape\raggedright \hspace*{10bp}\makebox[12bp][l]{\textbullet}ChatGPT uses deep learning, a subset of machine learning, to produce humanlike text through transformer neural networks. The transformer predicts text – including the next word, sentence or paragraph – based on its training data's typical sequence.\par}{\TimesNR\fontsize{24}{26.88}\selectfont\color{c000000}\itshape\vphantom{Ag}\par}{\TimesNR\fontsize{24}{26.88}\selectfont\color{c000000}\raggedright \hspace*{10bp}\makebox[12bp][l]{\textbullet}From \href{https://www.techtarget.com/whatis/definition/ChatGPT}{\nolinkurl{https://www.techtarget.com/whatis/definition/ChatGPT}} \par}{\TimesNR\fontsize{24}{26.88}\selectfont\color{c000000}\vphantom{Ag}\par}{\TimesNR\fontsize{24}{26.88}\selectfont\color{c000000}\raggedright \hspace*{10bp}\makebox[12bp][l]{\textbullet}Basically, it read a lot of documents and learnt their content as part of a bigger model of the knowledge in those documents. When generating answers, it starts ‘thinking with the moth open’ and determines each time the ‘most likely next word to say’. Under the hood its all about conditional probabilities!\par}\end{minipage}};
\node[anchor=north west,inner sep=0pt,outer sep=0pt] at (38.551,-65.764) {\begin{minipage}[t][73.709bp][c]{878.677bp}\setlength{\parindent}{0pt}\setlength{\parskip}{0pt}{\TimesNR\fontsize{44}{49.28}\selectfont\color{cBC0031}\bfseries\raggedright ChatGPT\par}\end{minipage}};
\end{abwframe}
